\section{Fine-Tuning Under an Information Budget}
\label{sec:fineqcomp}

\subsection{Related work and precise gap}
\TODO{Assume the reader already understands LLMs, LoRA, QLoRA, quantization, codelength, and rate--distortion from Section~\ref{sec:background}; do not re-explain them. Bring that well-versed reader to the specific modern literature on which this project builds. Cover only work that directly constrains or motivates our experiments: adaptive/structured PEFT where it bears on what an adapter can represent; quantization-aware and low-bit fine-tuning where it defines neighboring rate questions; post-training compression of trained LoRA adapters; MDL or prequential coding where it supplies operational notions of information; and recent work on adapter write/memorization capacity. For each neighboring line, make explicit what quantity it measures---trainable capacity, memory footprint, reconstruction error, memorization, or post-training artifact size---and why that differs from our quantity. End with the precise gap: for a fixed trained adapter, with the frozen base model available as side information, how many decoder-visible serialized bits are required to preserve its learned behavior, and can that rate be related to information in the training data relative to the receiver? This subsection should justify the experiments that follow rather than preview their results.}

\subsection{Starting point: rate--distortion pilots}
\TODO{Reconstruct the May \texttt{itft\_pilot} phase: rank sweeps, dataset-compressibility controls, 0.5B/1.5B confirmation, A/B asymmetry, and why ARC was rejected as the primary distortion measure.}

\figplaceholder{Early research figure: representative rank--held-out-loss curves and A-only/B-only/full comparison. Use only if they clarify the evolution of the question; otherwise move to appendix.}

\subsection{From nominal capacity to exact update rate}
\TODO{Explain why rank, nominal width, and raw parameter count were inadequate. Define exact serialized decoder-visible adapter rate and the fixed-checkpoint codec protocol.}

\figplaceholder{Pipeline diagram: dataset $\rightarrow$ QLoRA training $\rightarrow$ fixed raw adapter $\rightarrow$ codec ladder $\rightarrow$ serialized artifact $\rightarrow$ reload $\rightarrow$ behavioral evaluation.}

\subsection{Experimental protocol}
\TODO{Models, tasks, LoRA placement/rank, seeds, training protocol, exact-rate accounting, held-out NLL and exact-task evaluation. Clearly separate training from post-training compression.}

\subsection{Adaptive MDL allocation: hypothesis and negative result}
\TODO{Research sequence: hypothesize that reconstruction-aware allocation should preserve behavior efficiently; test adaptive allocation against fixed-rate controls; report the negative result.}

\figplaceholder{Headline rate--behavior frontier: adaptive MDL vs uniform / LoRAQuant-style controls.}

\subsection{Weight error does not determine behavioral error}
\TODO{Center the counterexample: adaptive allocation can obtain lower RMSE while catastrophically losing behavior relative to a slightly worse-RMSE fixed codec. State only after exact numbers are re-audited from result artifacts.}

\figplaceholder{Scatter: adapter reconstruction RMSE versus retained behavior, highlighting the decisive counterexample.}

\subsection{Failure analysis and experimental corrections}
\TODO{Row dropping catastrophe; exact-rate decomposition; cache contamination / accounting corrections; validation-opportunity bias; why these changed subsequent experimental design. Select only corrections with scientific relevance.}

\subsection{Does required update rate track dataset information?}
\TODO{Current compute-matched experiment: vary distinct information while holding optimizer updates and samples seen fixed; dense codec ladder around $\Rstar(0.9)$; compare unique rows, corpus code length, and behavioral $\Rstar$. Insert final results when campaign stabilizes.}

\figplaceholder{Primary ongoing plot: independently measured dataset information vs $\Rstar(0.9)$ with seed uncertainty.}

\subsection{Takeaways}
\TODO{Summarize what was learned about measuring fine-tuning information, including negative results and the methodological refinement they forced.}
